\documentclass[12pt, letterpaper]{article}
\date{\today}
\usepackage[margin=1in]{geometry}
\usepackage{amsmath}
\usepackage{hyperref}
\usepackage{cancel}
\usepackage{amssymb}
\usepackage{fancyhdr}
\usepackage{pgfplots}
\usepackage{booktabs}
\usepackage{pifont}
\usepackage{amsthm,latexsym,amsfonts,graphicx,epsfig,comment}
\pgfplotsset{compat=1.16}
\usepackage{xcolor}
\usepackage{tikz}
\usetikzlibrary{shapes.geometric}
\usetikzlibrary{arrows.meta,arrows}
\newcommand{\Z}{\mathbb{Z}}
\newcommand{\N}{\mathbb{N}}
\newcommand{\R}{\mathbb{R}}
\newcommand{\Q}{\mathbb{Q}}
\newcommand{\C}{\mathbb{C}}
\newcommand{\F}{\mathbb{F}}

\newcommand{\Po}{\mathcal{P}}
\newcommand{\Pro}{\mathbb{P}}
\author{Alex Valentino}
\title{501 homework}
\pagestyle{fancy}
\renewcommand{\headrulewidth}{0pt}
\renewcommand{\footrulewidth}{0pt}
\fancyhf{}
\rhead{
	Homework 3\\
	501
}
\lhead{
	Alex Valentino\\
}
\begin{document}
\begin{enumerate}
	\item[1] Let $\mathcal{M}$ be a $\sigma$-algebra and $(U_n){n \in \N} \subset \mathcal{M}$ be sequence contained in $\mathcal{M}$ with the property that all $U_n$ 
	are distinct.  We want to show that $\mathcal{M}$ is uncountable.  Consider the sequence 
	$C_1 = U_1, C_2 = U_2 \backslash U_1, C_m = U_m \backslash \bigcup_{i=1}^{m - 1} U_i$.  Note that since the sets are distinct then that implies there exists for 
	every $j \in \N$ a unique $x_j \in X$ such that $x_j \in U_j$.  Thus 
	$\emptyset \neq \{x_{j}\} \subset C_j$.  We claim that $C_j \cap C_i = \emptyset$
	for all $i,j \in \N, i \neq j$.  Assume WLOG that $i < j$.  Then $C_j$ explicily have any 
	elements in $U_i$.  Since $C_i \subseteq U_i$ then $C_i \cap C_j = \emptyset$.  
	Now since each $C_i$ is guarenteed to be unique, then we can construct the bijection
	between counable binary sequences and the countable union of elements in $(C_i)$.
	Note that this is achieved by assigning each binary digit, $b_i$, to whether 
	$C_i$ is included in the union.  Note that by the distinctness of each element in 
	$(C_i)$ that each union is unique.  Since the set of countable binary sequences is 
	uncountable, then the set of countable union of elements from the constructed sequence is countable.  Therefore since our constructed sequence is in $\mathcal{M}$
	then $\mathcal{M}$ necessarily contains an uncountable number of elements.  It can't
	have more than an uncountable number since only countable unions can be taken.  
	\item[4] Let $(X, \mathcal{M})$ be a measurable space and $f: X \to \R$.  We want 
	to show that if $A \cup B = X, A, B \in \mathcal{M}$, then $f$ is $\mathcal{M}, \mathcal{B}_\R$ measurable iff the restriction of $f$ to $A$ and $B$ is measurable too.
	\begin{itemize}
		\item $\Rightarrow$ Let $f$ be $\mathcal{M},\mathcal{B}_\R$ measurable.  
		Note that the restriction to $A$ or $B$ is just a subset of $\mathcal{M}$ with 
		the restriction that $f_A^{-1}(S) \subseteq A$.  If values that are outside 
		of the range of $f_A$ are asked of it then this still doesn' affect what $f_A^{-1}$ returns.  Thus for any $S \in \mathcal{B}_\R,$ $f_A^{-1}((a,b)) = f^{-1}((a,b)) \cap A \in \mathcal{A}$, where $\mathcal{A}$ is the algebra over $A$ with 
		respect to $\mathcal{M}$.  Note that the same applies to $B$. 
		\item $\Leftarrow$.  For all $S \in \mathcal{B}_\R$, 
		$f_A^{-1}(S) \in \mathcal{A},f_B^{-1}(S) \in \mathcal{B}$, $\mathcal{A}, \mathcal{B}$ are the algebras over $A,B$ contained in $\mathcal{M}$.  Note that 
		the inverse of the restriction is equivalent to taking the intersection of the function inverse and the chosen set.  Therefore $f^{-1}(S) = f^{-1}(S) \cap X = f^{-1}(S) \cap (A \cup B) =
		f^{-1}(S) \cap A \cup f^{-1}(S) \cap B = f_A^{-1}(S) \cup f_B^{-1}(S)$, which 
		both of those are in $\mathcal{M}$, thus $f^{-1}(S) \in \mathcal{M}$
	\end{itemize}
	\item[5] Let $f: \R \to \R$ be monotone.  Note that $f$
	is discontinuous in only countably many places.  Let the continuous portions (on their disjoint open set restriction)
	be denoted by the sequence $\{f_n\}_{n \in \N}$, allowing us to write 
	$f(x) = \cup_{i=1}^\infty f_i(x)$.    Let $S \in \mathcal{B}_\R$.  
	Consider $f^{-1}(S)$.  Note that by our definition $f^{-1}(x) = \bigcup_{i=1}^\infty f^{-1}_i(x)$.  Since each $f_i^{-1}$ is continuous on it's restriction, which is an open
	set in $\R$, then each $f_i^{-1}(S) \in \mathcal{B}_\R$, thus the countable 
	union of sets in the borel algebra is in the borel algebra.  Therefore 
	$f^{-1}(S) \in \mathcal{B}_\R$
	\item[7] Let $(X, \mathcal{M}, \mu)$ be a measure space with $\mu(X) < \infty$, $(f_n)_{n \in \N}$ be given by $f_i: \mathcal{M} \to \R$ as a sequence of measurable functions, and suppose there exists $E \in \mathcal{M}, \mu(E) = 0$, where 
	for all $x \in E^c, \sup_{n \in \N}|f_n(x)| < \infty$.  Consider the sets 
	$S_n = \{x: \sup_{n \in \N}|f_n(x)| > n\}$ where $n \in \N$.  Note that $S_{i+1} \subseteq S_{i}$ since it is possible that there exists elements $x \in S_i$ such that 
	$\sup_{n \in \N} |f_n(x)| < i + 1$.  Therefore $\lim \mu(S_n) = \mu(\lim S_n)$.
	Note that if $x \in X, \sup_{n \in \N}|f(x)| = \infty$ then necessarily, $x \in E$.
	Additionally this implies that $x \in \bigcap_{i=1}^\infty S_i$.  Since 
	$\bigcap_{n=1}^\infty S_n = \lim S_n$, then necessarily 
	$\lim S_n \subseteq E$.  Then $\mu(\lim S_n) \leq \mu(E) = 0$, thus $\lim \mu(S_n) = 0$, which implies for all $\epsilon > 0$ we can always find a $n$ sufficently large such
	 that $\mu(S_n) < \epsilon$.  Note that if $\mu(X) < \infty$ is relaxed then we 
	 can take $X = \N$ and have $\mu(S) = |S|$.  Note that the functions $f_n(x) = 
	 \begin{cases} x & x < n\\ n & x >= n \end{cases}$ are measurable functions since 
	 any open set in $\R$ containing a natural number between $1$ and $n$ will 
	 just include that in the inverse image of $f_n$, and since all finite subsets are 
	 in the $\sigma$ algebra of $\N$.
	 Additionally, for each $y \in \N$, $\sup_{n \in \N} |f_n(y)| = y$.  Therefore 
	 $E = \emptyset$.  However for any $C> 0$, $\mu (\{x \in \N: \sup_{n \in \N}|f_n(x)| > C\}) = \infty$.    
	\item[8] Let $(X, \mathcal{M}, \mu)$ be a measure space with $\mu(X) < \infty$.  
	Let $\{E_n\}_{n \in \N} \subset \mathcal{M}$, and let $E$ be the set of 
	all points such that $x \in E_n$ for an infinite number of $n$, and let $\mu(E) = 0$.  Let $F_1 = \bigcup_{n=1}^\infty E_n$, and $F_i = \bigcup_{n = i}^\infty E_n$.
	Note that $F_{i+1} \subseteq F_i$ and $\bigcap_{i=1}^\infty F_i = \bigcap_{i=1}^\infty E_i$.  Since $x \in \bigcap_{i=1}^\infty E_i$ implies that $x \in E$ then 
	$\bigcap_{i=1}^\infty E_i \subset E$.  Therefore $\lim \mu(E_n) = \mu(\bigcap_{n=1}^\infty E_n) \leq \mu(E) = 0$, which implies that $\lim \mu(E_n) = 0$.  If the condition 
	that $\mu(X) < \infty$ is relaxed then taking $X = \R$ we have the sequence 
	$E_n = (n,n+1)$.  Note that $E = \emptyset$, thus $\mu(E) = 0$, however 
	$\mu(E_n) = 1$.  
\end{enumerate}
\end{document}
